\scanpage{80}
\[
\underline H_Z^q(F)_{\bar x}\ne0\ \longrightarrow\ \dim\underline O_{Z,x}\deleted{\unsure{\leq}}\geq q-1
\]
i.e. $\underline H_Z^q(F)_{\bar x}$ est nul en codim $<q-1$,
[ce qui impliquera que $H^p(Z,\underline H_Z^q(F))\ne0\longrightarrow$
$p\leq(n-1)-(q-1)=n-q\deleted{\unsure{+1}}$ i.e. $p+q\leqn\deleted{\unsure{+1}}$, d'où
la conclusion voulue]. Il faut donc prouver que $q-1>\dim\underline O_{Z,x}\deleted{\unsure{+1}}\mathrel{\unsure{=|\Rightarrow}}\underline H_Z^q(F)_{\bar x}=0$.
Pour $q=0,1$, c'est \unsure{vide|évid|évidemment}. Sinon, on a
\[
\underline H_Z^q(F)_{\bar x}=H^{q-1}(X_{\bar x}-V(\pi)_{\bar x}\deleted{\unsure{d}},F),
\]
qui est \inserted{bien} nul
si $q-1>\dim X_{\bar x}$ [car \inserted{$X_{\bar x}$} c'est un
affine ess. de t.f. sur un corps sép. clos $\unsure{\bar K|K}$]
\deleted{de dim}
% slanted
\margin{\uline{\unsure{Calcul|Couac|Comme}} \uline{d'une} \uline{unité}}

\deleted{\uline{Question}} \inserted{\uline{Th}} Soit $A$ \uuline{str.} local, \inserted{$Y=\operatorname{Spec} A$, $Y'=Y-\{\unsure{m|\text{pt}}\}$,} $X$ \inserted{\unsure{aussi|ass.}} de type fini
et affine, \deleted{de dim de Krull} \inserted{$X'=X\times_Y Y'$}
Soit \deleted{\unsure{$d\leq$|$X_0$}} $d=\sup[\dim X_0,\ \dim X'+1]$.
Alors est-il vrai que pour tt $\ell$
premier à la car. rés. de $A$, on a
\[
\deleted{\unsure{H^i}}\ \operatorname{cd}_{\ell}(X)\leq d
\]
% slanted
\margin{Prouvé pour l'instant pour les anneaux ess. de type fini sur un anneau de dim $\leq1$}

\inserted{N.B. Si $X$ est intègre et domine $Y=\operatorname{Spec} A$ [ce qu'on peut tjrs supposer) on trouve $\operatorname{cd}_{\ell}(X)\leq\dim Y+d$ \inserted{$d=\dim$ de la fibre générique}}

N.B. Si on ne suppose pas $X$ affine, la formule serait
\[
\operatorname{cd}_{\ell}(X)\leq\sup[2\dim X_0,2\dim X'+1]
\]
\margin{\deleted{$\unsure{A\to Y}$,}}
resp. ($X$ irréd dominant) $\operatorname{cd}_{\ell}(X)\leq2\dim Y-1+2d$ \inserted{(si $\dim Y\ne0$)}
% archive sheet number: 79
